
\subsection{Reduced Error}


Errors can be thought of as assignments on the $d^{\prime}$-dimensional faces that have non-vanishing boundaries. In particular, $X$-errors can be characterized as $d^{\prime}$-dimensional chains $z \in \mathbb{F}^{\mathcal{F}_{d^{\prime}}}_{2}$ such that $\partial^{-} z \neq 0$. In this section, we assume dealing with the $\infty$-extension of a circuit (see \Cref{def:extension}), so the cubical complex is the $d$-dimensional grid with infinite side length. The \textit{reduced} form of an error is its representative, a translation by a stabilizer, which minimizes its size.  Here we prove that the sizes of a reduced error and its boundary (derived from $\partial^{-}$) are the same. In particular, they share poles of the potential walls. That will come into use in analyzing the behavior of the sweep-rule, as it is defined to decide on correction upon the syndrome it sees.

For doing so we present two codes that describe local views of vertices. They are the \textit{stabilizers code at vertex} $v$ and the \textit{small code at vertex} $v$. The first corresponds to stabilizers rooted at the vertex, while the second is the collection of local views on its positive faces that satisfy its positive checks. Then, in \Cref{claim:agreement}, we consider a vertex whose $X$-checks don't admit a syndrome, meaning that its local view is in its \textit{small code}, and show that there is a codeword of its \textit{stabilizers code} that agrees with its positive faces. Addition by that stabilizer yields a new chain excluding $v$. If $v$ were a unique maximum of $L_{d+1}$, then the vertices on the yielded chain would have a lower value of $L_{d+1}$. Subsequently, since the exclusion decreases the size, it follows that the 'corner'—the vertex achieving the maximal $L_{d+1}$ value—of a reduced error sees a syndrome.

\begin{definition}[Stabilizers Code at Vertex $v$]
We define \textit{stabilizer code at vertex} $v$, denoted by $\text{Stab}_{v}$, as the restriction of the $(d,d^{\prime})$-Toric code to the $d^{\prime}$-dimensional faces such that their vertices are at a distance of at most one positive step in each direction from $v$. That is equivalent to the union of all the positive $d^{\prime}$-dimensional faces of $v$ with the $d$-dimensional faces of each of its sibling $gv$, excluding faces containing $g$ as their generator. Namely:
  \begin{align*}
    \text{Stab}_{v} \colonequals  \bigl\{ z : \partial^{-} z  &= 0 \text{ and } \\
    & z \in \text{Span} \big(  \big\{  \rotfvs{} : S^{\prime} \subset S, |S^{\prime}| = d^{\prime}  \big\} \\
   & \ \ \ \ \ \ \ \ \ \ \ \ \  \bigcup_{g \in S} \big\{ \rotfv{gv}{S^{\prime}/g} : S^{\prime} \subset S/g , |S^{\prime}|=d^{\prime} \big\} \big) \bigr\}
  \end{align*}
\end{definition}
The checks of the code, are of the following three types: 
\begin{enumerate}
  \item Checks that are rooted at $v$, each identified by choosing $d-1$ generators from $S$. Namely: 
    \begin{equation*}
      \begin{split}
        \big\{ \rotfvs{}: S^{\prime} \subset S , |S^{\prime}|=d^{\prime}-1  \big\} 
      \end{split}
    \end{equation*}
    For each check, the bits that connect to it, are given by restricting the coboundary of the check to bits domain of $\text{Stab}_{v}$:  
    \begin{equation*}
      \begin{split}
        \big( \partial^{+} \rotfvs{} \big)|_{\text{Stab}_{v}} & =  \big(   \sum_{g \in S/S^{\prime}} \rotfv{v}{S^{\prime} \cup g} + \sum_{g \in S^{\prime}} \rotfv{g^{-1}v}{S^{\prime}\cup g}  \big)|_{\text{Stab}_{v}} =     \sum_{g \in S/S^{\prime}} \rotfv{v}{S^{\prime} \cup g} 
      \end{split}
    \end{equation*}
  \item Checks are rooted at the positive sibling of $v$. Consider such a sibling, let's say $gv$ (namely the vertex that is given by stepping in the direction $g \in S$ from $v$). Notice that any face rooted at $gv$ that contains $ggv$ cannot be contained in faces rooted at $v$. Thus, the checks of $gv$ correspond to choosing subsets from $S/g$. Namely:
    \begin{equation*}
      \begin{split}
      \bigcup_{g \in S}\big\{ \rotfv{gv}{S^{\prime}}: S^{\prime} \subset S/g , |S^{\prime}|=d^{\prime}-1  \big\}
      \end{split}
    \end{equation*}
    And each check is supported by:
    \begin{equation*}
      \begin{split}
        \big( \partial^{+} \rotfv{gv}{S^{\prime}} \big)|_{\text{Stab}_{v}} & = 
        \big(   \sum_{g^{\prime} \in S/S^{\prime}} \rotfv{gv}{S^{\prime} \cup g^{\prime}} + \sum_{g^{\prime} \in S^{\prime}} \rotfv{g^{\prime -1}gv}{S^{\prime}\cup g^{\prime}} \big)|_{\text{Stab}_{v}} \\
        &  =    \bigl(\sum_{ \substack{ g^{\prime} \in S/S^{\prime} \\ g^{\prime } \neq g }} \rotfv{g v}{S^{\prime} \cup g^{\prime}} \bigr) + \rotfv{v}{S^{\prime}\cup g}
      \end{split}
    \end{equation*}
  \item Checks that their roots are at a two-step distance from $v$, namely vertices of the form $g^{\prime}gv$ where $g, g^{\prime} \in S$, and $g \neq g^{\prime}$. By the arguments presented in the second type, the checks are the $d^{\prime}-1$ faces which do not have the generators on the path leading from $v$ to them, i.e., $g, g^{\prime}$, in their generator set.
    \begin{equation*}
      \begin{split}
        \bigcup_{g,g^{\prime} in S^{\prime} }\big\{ \rotfv{g^{\prime}gv}{S^{\prime}}: S^{\prime} \subset S/\{ g, g^{\prime} \} , |S^{\prime}|=d^{\prime}-1  \big\}
      \end{split}
    \end{equation*}
  And for each check, its support is:
    \begin{equation*}
      \begin{split}
        \big( \partial^{+} \rotfv{g^{\prime}gv}{S^{\prime}}  \big)|_{\text{Stab}_{v}} & =  \big(    \sum_{g^{\prime \prime} \in S/S^{\prime}} \rotfv{gg^{\prime}v}{S^{\prime} \cup g^{\prime \prime}} + \sum_{g \in S^{\prime}} \rotfv{g^{\prime \prime -1}gg^{\prime}v}{S^{\prime}\cup g^{\prime \prime}} \big)|_{\text{Stab}_{v}}\\
        &=    \rotfv{gv}{S^{\prime} \cup g^{\prime}} + \rotfv{g^{\prime}v}{S^{\prime}\cup g}   
      \end{split}
    \end{equation*}
\end{enumerate}

\begin{definition}[Small Code at Vertex $v$]
For a vertex $v$, we denote by $\mathcal{C}_{v}$ the \textit{small code at vertex $v$}, which consists of all the $d^{\prime}$-dimensional chains rooted at $v$ that satisfy the positive checks of $v$, namely:
  \begin{align*}
    \mathcal{C}_{v} \colonequals  \bigl\{ z : (\partial^{-} z)|_{v+}  &= 0 \text{ and } \\
    & z \in \text{Span} \big(  \big\{  \rotfvs{} : S^{\prime} \subset S, |S^{\prime}| = d^{\prime}  \big\} \big) \bigr\} 
  \end{align*}
The checks of the code are exactly the first type checks of $\text{Stab}_{v}$.
\end{definition}
\begin{claim}
  \label{claim:agreement}
  Let $v$ be a vertex. For any $d^{\prime}\ge 2$, $d\ge d^{\prime}+1$, and any $z \in \mathcal{C}_{v}$, there is a stabilizer $w \in \text{Stab}_{v}$ such that $z = w|_{v}$; namely, $z$ and $w$ agree on the faces rooted at $v$. Moreover, the vertex support of $w|_{v}$ is contained in the vertex support of $z|_{v}$, namely:
  \begin{equation*}
    \begin{split}
        \textbf{Str}_{0+}(v, w) \subset \textbf{Str}_{0+}(v, z)
    \end{split}
  \end{equation*}
\end{claim}

\begin{proof}
  By induction on the dimensions $d, d'$. We extend the notation of the stabilizer and the small codes at vertex $v$ to $\mathcal{C}_{v}^{J}, \text{Stab}_{v}^{J}$, where $J$ is a subset of $S$, to refer to the codes derived from the skeleton generated by $J$. Showing that the belonging $z\in \mathcal{C}_{v}^{J}$ implies the existence of $w \in \text{Stab}_{v}^{J}$ such that $w|_{v} =z|_{v}$ proves the claim.
  \begin{enumerate}
    \item Base. For the base cases, let $d$ be either three or four, and let $d^{\prime}$ be two. In the appendix, we show that the generator matrix of $\mathcal{C}_{v}$ is a submatrix of the generating matrix of $\text{Stab}_{v}$. See \Cref{subsection:basecase}.
      
    \item Assume the correction of the claim for dimensions lower than $d, d^{\prime}$. Decompose $z = z|_{gv} + z|_{/gv}$ where $z|_{gv}$ and $z|_{/gv}$ are the chains consisting of all the faces in $z$ that have support on $gv$ and its complement. Define $z^{*}$ to be the $d^{\prime}-1$-dimensional chain obtained by excluding $gv$ from any face of $z|_{gv}$. Notice that $\partial^{-} z^{*}$ has to be zero since it contains all the faces of $\partial^{-}z$ that contain $gv$ (up to the removal). In other words, $z^{*} \in \mathcal{C}_{v}^{S/g}$. Thus, we can use the induction assumption and infer the existence of $w^{*} \in \text{Stab}_{v}^{S^{\prime}/g}$ such that $z^{*} = w^{*}|_{v}$.

      
Now, since $w^{*} \in \text{Stab}_{v}^{S/g}$, then $\partial^{-}w = 0$. In particular, $w^{*}$ can be decomposed into a sum of stabilizers rooted in $v$. Let $\rotfv{v}{S^{\prime}_{1}}, \rotfv{v}{S^{\prime}_{2}}, \rotfv{v}{S^{\prime}_{3}}, \ldots, \rotfv{v}{S^{\prime}_{l}}$ be the $d^{\prime}$-dimensional chains corresponding to the stabilizers composing $w^{*}$. Namely:
      \begin{equation*}
        \begin{split}
          w^{*} = \sum_{i} \partial^{-} \rotfv{v}{S^{\prime}_{i}} 
        \end{split}
      \end{equation*}
Since those stabilizers are basis elements of $\text{Stab}_{v}^{S/g}$, we have that $g \notin S^{\prime}_{i}$ for any $i$. Hence, we can define the extended $d^{\prime}$-dimensional chain $w$ to be the chain obtained by adding $g$ to each of the stabilizers deriving $w^{*}$. Namely:
      \begin{equation*}
        \begin{split}
          w =  \sum_{i} \partial^{-} \rotfv{v}{S^{\prime}_{i} \cup g} 
        \end{split}
      \end{equation*}
However, $\rotfv{v}{S^{\prime}_{1} \cup g}, \ldots, \rotfv{v}{S^{\prime}_{l} \cup g}$ are $d^{\prime}+1$-dimensional chains corresponding to the stabilizers, which are basis elements of $\text{Stab}_{v}^{S} = \text{Stab}_{v}$. Therefore, $w \in \text{Stab}_{v}$. In addition, opening up the boundary operator yields:
      \begin{equation*}
        \begin{split}
          w|_{v} &=  \sum_{i} \bigl(\partial^{-} \rotfv{v}{S^{\prime}_{i} \cup g} \bigr)|_{v} = \sum_{i} \bigl(\sum_{g^{\prime} \in S^{\prime}_{i}/g}\rotfv{v}{S^{\prime}_{i}/g^{\prime} \cup g} + \rotfv{v}{S^{\prime}_{i}}   \bigr)\\
          &= \sum_{\rotfvs{}\in w^{*}|_{v} }\rotfv{v}{S^{\prime} \cup g} + \sum_{i} \rotfv{v}{S^{\prime}_{i}} = z|_{gv} + \sum_{i} \rotfv{v}{S^{\prime}_{i}} 
        \end{split}
      \end{equation*}
Thus, adding $w|_{v}$ to $z$ results in:
      \begin{equation*}
        \begin{split}
          z + w|_{v} = z|_{/gv} + \sum_{i} \rotfv{v}{S^{\prime}_{i}} 
        \end{split}
      \end{equation*}
Which is a $d^{\prime}$-dimensional chain, supported on vertices of the form $g^{\prime}v$ for $g^{\prime}\in S/g$, namely a $d^\prime$-dimensional chain on the complex induced by the $d-1$-dimensional skeleton. In addition, since the boundary operator can only advance the root of a face forward, we have that $\partial^{-} w|_{v} =  (\partial^{-} w)|_{v}$ and therefore $\partial^{-}( z + w|_{v}) = \partial^{-}z + (\partial^{-}w)|_{v}$.
      
To summarize, we have both $z + w|_{v} \in \mathcal{C}_{v}^{S/g}$ and $z+w|_{v}$ sets on a $d-1$-dimensional skeleton. Therefore, we can repeat in a recursive manner, with each step resulting in a $d^{\prime}$-dimensional chain in a skeleton with dimension decreased by one. We stop when $d = d^{\prime}$; in that case, the resulting chain has to be trivial.
  \end{enumerate}
\end{proof}

\begin{claim}
  \label{claim:reducederr}
  Let $x \in \mathbb{F}_{2}^{\mathcal{F}_{d^{\prime}}}$ be a $d^{\prime}$-dimensional finite chain, consider the stabilizer  $z \in \mathcal{C}_{Z}^{\perp}$ which adding it to $x$ minimize the size , Namely: 
  \begin{equation*}
    \begin{split}
      z =  \arg \min_{z \in \mathcal{C}_{Z}^{\perp}} \Size{ x + z }  
    \end{split}
  \end{equation*}
  Then: 
  \begin{equation*}
    \begin{split}
      \Size{ x + z } = \Size{ \partial^{-}( x + z) } 
    \end{split}
  \end{equation*}
\end{claim}
\begin{proof}
  First, observe that since the vertex domain of $\partial^{-} x$ is contained in the vertex domain of $x$, it follows that for any $x$ the inequality $\Size{x} \ge \Size{\partial^{-} x}$ holds. Combining this with the fact that $\partial^{-} z = 0$ for any $z \in \mathcal{C}_{Z}^{\perp}$, we have that $\Size{x + z} \ge \Size{\partial^{-}(x + z)}$. Now suppose that $\Size{x} > \Size{\partial^{-} x}$; that implies the existence of $v \in x/\partial^{-} x$ such that $v$ is a local maximum of one of the potential walls in $x$ and $(\partial^{-} x)|_{v} = 0$. Let's split into cases, depending on the potential wall type.
 
  \begin{enumerate}

      \item If $v$ is a maximum of $L_{d+1}$, then by \Cref{claim:agreement}, there exists a stabilizer $z$ rooted at $v$, such that $z|_{v}=z|_{v+}=x|_{v+}=x|_{v}$ (where we used that $v$ is a local maximum of $L_{d+1}$ in $x$ to infer $x|_{v}=x|_{v+}$). Therefore, $(x+z)|_{v} = 0$, meaning adding $z$ to $x$ excludes $v$. In addition, since the rooted stabilizer $z$ is supported only on the vertices in $\textbf{Str}_{0+}(v,x)$, the addition cannot introduce a vertex with a higher $L_{g}$ value than the maximal $L_{g}$ value in $\textbf{Str}_{0+}(v,x)$ into $x$. So, repeating the process until excluding all the vertices that get the maximal $L_{d+1}$-value in $x$, and summing the stabilizers, results in a stabilizer $z$ for which $L_{d+1}(x+z) < L_{d+1}(x)$.


      \item If, for some generator $g \in S$, $v$ is a maximum of $L_{g}$, then without loss of generality, pick $v$ among the maximum points of $L_{g}$ with the highest value of $-L_{d+1}$ (maximize $\sum_{g \in S} L_{g}(v)$). Since stepping through the positive direction can only increase the value of $-L_{d+1}$, it follows that $v$ is also a local max of $-L_{d+1}$. Now consider the skeleton system where the directions are flipped, and $\rotfs{}$ is set to be the face containing the vertices that are reached by the selection of reverse steps from $S^{\prime}$ (i.e., $g^{-1}v \in \rotfs{}$). In that system, we can apply \Cref{claim:agreement} (in the reversed direction) and get that there is a stabilizer $z$ such that adding it to $x$ excludes $v$ from $x$. By the same argument as above, no vertex with a higher $L_{g^{\prime}}$-value (although $L_{d+1}$ can increase, that's fine) than what it already has in $\textbf{Str}_{0-}(v,x)$ is introduced into $x$. Repeating the process excludes all the vertices with the maximal value of $L_{g}$ in $x$, so eventually, we end up with $x+z$, such that $L_{g}(x+z) < L_{g}(x)$.
  \end{enumerate}

  
We keep excluding vertices until we get the desired equality.

\end{proof}

We close the section by defining the \textit{reduced} form of a finite $d^{\prime}$-dimensional chain and the notion of \textit{finite} error.
\begin{definition}
  \label{def:reduced}
  Let $x \in \mathbb{F}_{2}^{\mathcal{F}_{d^{\prime}}}$ be a $d^{\prime}$-dimensional finite chain. We will say that $x$ is reduced if for any vertex $v \in x$ that is a maximum point of a potential wall $L \in L_{1}, \ldots, L_{d+1}$, $v$ is also in $\partial x$. 
In particular, we have that $\Size{x} = \Size{\partial x}$ and that the reduced form of a given chain is obtained by the construction presented in the proof of \Cref{claim:reducederr}.
\end{definition}

\begin{figure}[h]
  \centering
  \begin{tikzpicture}[scale=1.5]

\tikzset{xyp/.style={canvas is xy plane at z=#1}}
\tikzset{xzp/.style={canvas is xz plane at y=#1}}
\tikzset{yzp/.style={canvas is yz plane at x=#1}}



\draw[blue!20!black,fill=blue!70!black!70,rounded corners=1,opacity=0.5, xzp=0]%,shift={(0.15,0.15)}]
(-0.1,2.1) rectangle (1.1,0.9);


      \draw[blue!20!black,fill=blue!70!black!70,rounded corners=1,opacity=0.5, yzp = 0]%,shift={(0.15,0.15)}]
      (-0.1,2.1) rectangle (1.1, 0.9);

      \draw[blue!20!black,fill=blue!70!black!70,rounded corners=1,opacity=0.5, xyp = 2]%,shift={(0.15,0.15)}]
      (-0.1,-0.1) rectangle (1.1, 1.1);
      
      \draw[blue!20!black,fill=blue!70!black!70,rounded corners=1,opacity=0.5, yzp = 1]%,shift={(0.15,0.15)}]
      (-0.1,2.1) rectangle (1.1, 0.9);
%,opacity=0.5]
%\fill[blue!30]
%(0.8,0.8,1.2) -- (0.7,0.7,0) -- (2,2,0) -- (2,2,1.2) -- cycle;
%
%\fill[blue!30]
%(0.9,0.9,1.1) -- (0.9,0.9,-0.1) -- (1.1,0.9,-0.1) -- (1.1,0.9,1.1) -- cycle;  

      \node[below = 3pt] at (0,0,2) { $v$ } ;

% draw grid lines
\foreach \x in {0,1,2} {
  \foreach \y in {0,1,2} {
    \draw[gray] (\x,\y,0) -- (\x,\y,2);
  }
}

\foreach \x in {0,1,2} {
  \foreach \z in {0,1,2} {
    \draw[gray] (\x,0,\z) -- (\x,2,\z);
  }
}

\foreach \y in {0,1,2} {
  \foreach \z in {0,1,2} {
    \draw[gray] (0,\y,\z) -- (2,\y,\z);
  }
}

% draw lattice points
\foreach \x in {0,1,2} {
  \foreach \y in {0,1,2} {
    \foreach \z in {0,1,2} {
      \fill (\x,\y,\z) circle (1pt);
    }
  }
}

\draw (2.5,1) node { $\rightarrow$ };
\end{tikzpicture} \begin{tikzpicture}[scale=1.5]

\tikzset{xyp/.style={canvas is xy plane at z=#1}}
\tikzset{xzp/.style={canvas is xz plane at y=#1}}
\tikzset{yzp/.style={canvas is yz plane at x=#1}}

\draw[green!20!black,fill=green!70!black!70,rounded corners=1,opacity=0.5, xzp=1]%,shift={(0.15,0.15)}]
(-0.1,2.1) rectangle (1.1,0.9);

      \draw[green!20!black,fill=green!70!black!70,rounded corners=1,opacity=0.5, xyp = 1]%,shift={(0.15,0.15)}]
      (-0.1,-0.1) rectangle (1.1, 1.1);

      \node[below = 3pt] at (0,0,2) { $v$  } ;
      \node[below] at (0,1,2) { $u$ } ;
      \node[below] at (1,0,1) { $w$ } ;

% draw grid lines
\foreach \x in {0,1,2} {
  \foreach \y in {0,1,2} {
    \draw[gray] (\x,\y,0) -- (\x,\y,2);
  }
}

\foreach \x in {0,1,2} {
  \foreach \z in {0,1,2} {
    \draw[gray] (\x,0,\z) -- (\x,2,\z);
  }
}

\foreach \y in {0,1,2} {
  \foreach \z in {0,1,2} {
    \draw[gray] (0,\y,\z) -- (2,\y,\z);
  }
}

% draw lattice points
\foreach \x in {0,1,2} {
  \foreach \y in {0,1,2} {
    \foreach \z in {0,1,2} {
      \fill (\x,\y,\z) circle (1pt);
    }
  }
}
\end{tikzpicture}
\caption{
Consider the $(3,2)$-Toric code, with qubits and $X$-checks set on the 2- and 1-dimensional faces. In blue, we have a 2-dimensional chain where the vertex achieving the minimum of $L_{d+1}$, denoted by $v$, does not see any unsatisfied checks in its positive local view. The green faces represent a 2-dimensional chain, which is the \textbf{reduced} form of the blue. In particular, any local maximum of $L_{d+1}$ (vertices $u$ and $w$) sees a non-trivial syndrome in its positive local view. Namely, the maximal points of $L_{d+1}$ in the chain are also the maximal points in the boundary $\partial^{-}$.
}
  \label{fig:edge}
\end{figure}



\begin{definition}
  \label{def:finiteerr}
Let $z$ be a $d^{\prime}$-dimensional chain. Consider the graph whose vertices are the faces of $z$, and any two vertices are connected by an edge if their intersection is a $d^{\prime}\pm 1$ dimensional face. Then $z$ is said to be connected if that graph is connected. In addition, we will say that a connected chain $z$ is finite if $\Size{z} < n$. Similarly, $z$ is infinite if $\Size{z} \ge n$. We extend the finite notation to unconnected chains and Pauli operators as follows: an unconnected chain will be said to be finite if any subset of it which is connected is finite. A Pauli operator will be said to be finite if the support of its view in the $\infty$-extension is a finite chain.
\end{definition}

